\section{Discussion}\label{sec:discussion}

Our 21-year survey revealed substantial spatial and temporal variation in esca incidence across Western Switzerland, with climatic factors related to water balance emerging as important drivers at the regional scale. This network was developed as a surveillance project for regional risk assessment and covers the main viticultural areas of western Switzerland. We analysed the data at two spatial scales: the vineyard level and the regional level.

Over the years of monitoring, esca incidence varied from year to year, both within and across vineyards. This is consistent with the fluctuations in esca incidence reported between years and plots in other long-term monitoring studies \citep{dewasme_long-term_2022,etienne_exploring_2024,etienne_identification_2026}. However, we observed a decrease in the overall incidence of esca, contrary to other studies reporting an increase \citep{bruez_overview_2012} or a stable trend \citep{etienne_exploring_2024}. Higher temperature and lower precipitation are considered factors that decrease esca incidence \citep{bortolami_grapevines_2021,etienne_exploring_2024} whereas higher precipitation in early summer is positively correlated with esca incidence \citep{calzarano_rainfall_2018,serra_expression_2018}. \citep{scherrer_trends_2022} The increasingly dry conditions observed in Western Switzerland in past two decades, particularly during summer, resulting from changes in precipitation and increased evaporative demand associated with rising temperatures (Scherrer et al., 2022), may have contributed to the observed decrease in esca incidence. These findings highlight the value of long-term monitoring programmes for tracking disease dynamics under changing climatic conditions.

While a decreasing trend was observed across the network, substantial differences emerged among regions. In La Côte, Chablais, Neuchâtel and Lavaux, where esca incidence was reported to be higher, the decreasing trend was more pronounced. Valais and Nord Vaudois, which have historically lower rates of esca incidence, showed a more stable long-term trend. These regional disparities likely reflect differences in climatic conditions and soil characteristics \citep{calzarano_rainfall_2018,marchi_observations_2006,monod_identifying_2025}. Valais is characterized by a drier climate \citep{zumbrunnen_linking_2009} and well-drained soils \citep{sigales_etude_2007} compared to the other regions of our network, which can contribute to explain its lower and more stable esca incidence. In Nord Vaudois, the parent material lies closer to the surface, giving rise to soils with reduced water-retention capacity. Drier climate and poor soil water-retention capacity are therefore factors that reduce esca incidence \citep{bortolami_grapevines_2021,etienne_identification_2026,monod_identifying_2025}. Interestingly, substantial variation was also observed among plots within the same wine-growing region despite comparable climatic conditions, suggesting that local factors may play a critical role in disease expression.

We assessed the predictive power of several models using meteorological variables to predict esca incidence (\% of symptomatic plants per plot) across our plot network, but the performance remained poor at the plot level. The low prevalence of symptomatic vines (approximately 2-3\%) combined with relatively limited interannual variation reduced the signal available to modelling approaches, making accurate predictions inherently challenging. Esca incidence cannot be explained by meteorological variables alone, in line with previous studies conducted in various geographical contexts \citep{bruez_overview_2012,dewasme_long-term_2022,etienne_identification_2026,lecomte_logistic_2024}. These findings suggest that intrinsic plot characteristics may contribute more strongly to disease expression than meteorological variables alone, suggesting that other, still unidentified factors are also involved.

The output of the FORD algorithm shows that, at the monthly scale, VPD was the variable most consistently associated with esca incidence, whereas at the seasonal scale, precipitation and temperature showed the strongest associations. Among all variables tested, VPD also appeared to be the most stable predictor across models, which is known to be an important driver ~of plant functioning \citep{grossiord_plant_2020}. Although the overall predictive performance of the models remained modest, these findings are consistent with previous studies. For example, Etienne et al. (2026) reported a positive relationship between temperature and precipitation during the symptom expression period and a negative relationship between VPD during the growing season and esca incidence. Similarly, VPD, temperature, and precipitation have been identified as relevant climatic drivers of esca expression in other viticultural contexts, lending further support to our results. These variables are closely linked to seasonal dryness, which is increasingly recognized as a major determinant of esca incidence \citep{bortolami_grapevines_2021}.

Under drier conditions, elevated VPD increases atmospheric water demand and promotes stomatal closure, thereby reducing transpirational water loss \citep{grossiord_plant_2020}. One possible explanation is that reduced xylem transport under water-limited conditions could limit the movement of phytotoxic compounds involved in symptom expression, although experimental evidence supporting this mechanism remains limited \citep{bortolami_grapevines_2021}. Conversely, greater water availability may facilitate transpiration and xylem transport processes, potentially increasing the movement of such compounds within the plant.

Since our models did not achieve strong predictive power for esca incidence based on meteorological variables at the plot level, we performed regressions using the most promising variables (VPD and the precipitation-to-evapotranspiration ratio) against esca incidence by region. At the regional scale, a higher precipitation-to-evapotranspiration ratio was associated with an increase in esca incidence, which in turn are generally associated with greater vine vigor. In addition, higher precipitation is likely associated with increased soil water availability, reflected by a higher Soil Water Index (SWI) which is a determinant factor known to increase esca incidence \citep{etienne_identification_2026,monod_identifying_2025}. At the other end of the gradient, Valais, the driest region considered, consistently showed the lowest esca incidence. The correlation between high VPD and lower esca incidence across regions confirms this pattern: when atmospheric demand is higher, meaning a greater water deficit, esca incidence is lower. Since the regions covered by our dataset exhibit a wide range of hydrological regimes, our data set provide a reasonably uniform coverage of the interval of both ET/P and VPD variables. This regional-scale pattern is consistent with the variable importance and FORD results: VPD and water-related variables are the dominant climatic drivers of Esca expression. The relative concordance between the plot-level statistical analyses and this simple regional characterization reinforces the conclusion that water stress (or lack thereof) is the primary climatic factor modulating Esca symptom expression that our data can identify. This strengthens our confidence in the association between these variables and esca incidence. The growing body of evidence linking water availability with esca incidence suggests that irrigation practices should be carefully considered, particularly in susceptible cultivars. \citep{etienne_identification_2026}.

One of the key findings of this study is that meteorological variables explained variation in esca incidence more effectively at the regional than at the plot scale. This suggests that climate may primarily influence esca incidence at a broader spatial scale while local environmental, pedological, and biological factors may play a greater role in determining disease expression at the vineyard level.
